\chapter{Objetivos}
\label{cap:objetivos}

El objetivo general del presente Trabajo de Fin de Grado es el diseño, implementación y validación de un sistema de monitorización y control de confort térmico para el hogar basado en \textit{Edge AI} que opere de forma completamente autónoma en una red local privada, eliminando cualquier dependencia de servicios en la nube para su funcionamiento.

\section{Objetivos Específicos}

A continuación se desglosa cada objetivo específico en subobjetivos medibles que permitan verificar su cumplimiento:

\begin{enumerate}
    \item \textbf{Diseñar una arquitectura de microservicios local.}
    \begin{itemize}
        \item Definir las interfaces de comunicación entre cada componente del sistema (orquestador, módulo HVAC, módulo IoT, sensores), garantizando que sean independientes y sustituibles.
        \item Implementar el orquestador OpenClaw como servicio systemd con auto-arranque y reinicio automático.
        \item Documentar la arquitectura mediante diagramas de despliegue y de flujo de datos.
        \item Verificar que ningún componente requiere acceso a Internet para su funcionamiento básico.
    \end{itemize}

    \item \textbf{Implementar el control de la unidad HVAC por bus LIN.}
    \begin{itemize}
        \item Diseñar el firmware en C++20 para el ESP32 DevKit V1 que actúe como puente entre la red local (serie/USB) y el bus LIN de la máquina General/Fujitsu.
        \item Implementar el parser de tramas LIN basado en la documentación del proyecto \texttt{esphome-fujitsu-halcyon}~\cite{omniflux_esphome_fujitsu}.
        \item Validar la comunicación bidireccional: lectura de estado del termostato y envío de comandos de control.
        \item Garantizar que el firmware mantenga el último estado válido en caso de pérdida de comunicación con el servidor (RNF-04).
    \end{itemize}

    \item \textbf{Integrar dispositivos IoT heterogéneos.}
    \begin{itemize}
        \item Lectura de sensores Xiaomi LYWSD03MMC mediante BLE GATT directo, utilizando exclusivamente la librería \texttt{bleak}~\cite{bleak}, sin depender de la nube de Xiaomi ni de firmware alternativo (RF-02).
        \item Control local del enchufe inteligente Tuya mediante \texttt{tinytuya}~\cite{tinytuya} sobre la red air-gapped, sin pasar por los servidores de Tuya (RF-03). Tiempo de respuesta inferior a 2~s.
        \item Ingeniería inversa completa del protocolo BLE del ventilador FanLamp F8808, incluyendo el mapeo de todos los comandos funcionales (RF-04).
        \item Documentar cada integración con los pasos necesarios para reproducirla en otro entorno.
    \end{itemize}
\end{enumerate}

Todas las integraciones deben cumplir los siguientes criterios transversales:
\begin{itemize}
    \item \textbf{Coste cero en licencias o suscripciones:} no se utilizarán servicios de pago, APIs externas ni plataformas que requieran cuota.
    \item \textbf{Reproducibilidad:} cualquier persona con el mismo hardware debe poder replicar la integración siguiendo la documentación.
    \item \textbf{Validación sobre hardware real:} todas las pruebas deben realizarse sobre los dispositivos físicos en la red air-gapped, no en simulación.
\end{itemize}

\section{Objetivos de Evaluación}

Para medir el grado de cumplimiento del proyecto, se establecen los siguientes indicadores:

\begin{enumerate}
    \item \textbf{Cobertura del enrutador de comandos:} el sistema debe ser capaz de interpretar y responder correctamente al menos el 90\% de los comandos en lenguaje natural que recibe, sin requerir reformulación por parte del usuario.
    \item \textbf{Latencia de respuesta:} la mediana de latencia para comandos directos (encender/apagar, cambio de velocidad) debe ser inferior a 2~s en la red local.
    \item \textbf{Disponibilidad del sistema:} los servicios críticos (orquestador, sensor daemon, router de modelos) deben mantener una disponibilidad superior al 99\% en condiciones normales de operación.
    \item \textbf{Aislamiento de red:} debe verificarse que ningún paquete de datos del ecosistema abandone la subred \texttt{192.168.1.0/24} durante su funcionamiento normal.
\end{enumerate}

\section{Restricciones Asumidas}

El proyecto asume las siguientes restricciones que acotan el alcance de los objetivos:

\begin{itemize}
    \item El desarrollo del firmware HVAC (ESP32 + LIN) queda como objetivo principal pero con posible desplazamiento a trabajo futuro si las pruebas de integración no se completan dentro del calendario académico.
    \item No se implementará una aplicación móvil nativa; el control se realiza exclusivamente mediante Telegram y el \textit{dashboard} web de OpenClaw.
    \item La base de datos SQLite de analítica energética se considera un objetivo secundario, priorizando la funcionalidad de control en tiempo real.
\end{itemize}
